\documentclass[11pt]{article}

\usepackage[margin=0.9in]{geometry}
\usepackage{amsmath,amssymb}
\usepackage{booktabs}
\usepackage{enumitem}
\usepackage[hidelinks]{hyperref}

\setlist[itemize]{leftmargin=1.4em}
\setlist[enumerate]{leftmargin=1.6em}

\title{SYSC 5108 Exam Study\\Assignment 4, Question 1}
\author{Jesse Levine}
\date{}

\begin{document}
\maketitle

\section*{Question Summary}

The network is a simple recurrent neural network with:
\begin{itemize}
    \item two input neurons,
    \item three hidden ReLU neurons,
    \item two output ReLU neurons,
    \item all bias terms equal to \(0.1\).
\end{itemize}

Given
\[
    x_t = \begin{bmatrix}1 \\ 0.5\end{bmatrix},
    \qquad
    h_{t-1} = \begin{bmatrix}0.05 \\ 0.2 \\ 0.15\end{bmatrix},
\]
the question asks for:
\begin{enumerate}[label=(\alph*)]
    \item the output \(y_t\),
    \item the \(\delta\) value for each neuron when the target is
    \(\begin{bmatrix}0 \\ 1\end{bmatrix}\).
\end{enumerate}

\section*{Course and Textbook Cross-References}

\begin{itemize}
    \item \textbf{Chapter 4 slides, Forward Propagation Generalization:} each
    layer computes an affine transformation followed by an activation.
    \item \textbf{Chapter 4 slides, Hidden Units:} ReLU is
    \[
        g(z)=\max(0,z).
    \]
    \item \textbf{Chapter 4 slides, Chain Rule and Backward Propagation:}
    gradients propagate backward through the network using the chain rule.
    \item \textbf{Goodfellow, Bengio, and Courville, \emph{Deep Learning}:}
    feedforward layers use affine transformations plus activations, while recurrent
    networks extend this by feeding a previous hidden state into the current hidden
    update. ReLU units remain elementwise \(g(z)=\max(0,z)\).
\end{itemize}

\section*{Given Parameters}

Bias vectors:
\[
    b_h=\begin{bmatrix}0.1\\0.1\\0.1\end{bmatrix},
    \qquad
    b_y=\begin{bmatrix}0.1\\0.1\end{bmatrix}.
\]

Weight matrices:
\[
    W_{xh}
    =
    \begin{bmatrix}
        -0.07 & 0.05 \\
        -0.04 & 0.10 \\
        -0.05 & -0.05
    \end{bmatrix},
\]
\[
    W_{hh}
    =
    \begin{bmatrix}
        -0.22 & -0.10 & 0.05 \\
        -0.04 & -0.09 & -0.06 \\
        -0.09 & 0 & -0.08
    \end{bmatrix},
\]
\[
    W_{hy}
    =
    \begin{bmatrix}
        0.06 & -0.01 & -0.18 \\
        -0.06 & 0.13 & 0.14
    \end{bmatrix}.
\]

\section*{Step-by-Step Solution}

\subsection*{Part (a): Compute \(y_t\)}

For one recurrent step, the hidden pre-activation is
\[
    z_h = W_{xh}x_t + W_{hh}h_{t-1} + b_h.
\]

\subsubsection*{1. Compute \(W_{xh}x_t\)}

\[
\begin{aligned}
    W_{xh}x_t
    &=
    \begin{bmatrix}
        -0.07 & 0.05 \\
        -0.04 & 0.10 \\
        -0.05 & -0.05
    \end{bmatrix}
    \begin{bmatrix}1 \\ 0.5\end{bmatrix} \\
    &=
    \begin{bmatrix}
        -0.07(1)+0.05(0.5) \\
        -0.04(1)+0.10(0.5) \\
        -0.05(1)-0.05(0.5)
    \end{bmatrix} \\
    &=
    \begin{bmatrix}
        -0.045 \\
        0.010 \\
        -0.075
    \end{bmatrix}.
\end{aligned}
\]

\subsubsection*{2. Compute \(W_{hh}h_{t-1}\)}

\[
\begin{aligned}
    W_{hh}h_{t-1}
    &=
    \begin{bmatrix}
        -0.22 & -0.10 & 0.05 \\
        -0.04 & -0.09 & -0.06 \\
        -0.09 & 0 & -0.08
    \end{bmatrix}
    \begin{bmatrix}
        0.05 \\ 0.2 \\ 0.15
    \end{bmatrix} \\
    &=
    \begin{bmatrix}
        -0.22(0.05)-0.10(0.2)+0.05(0.15) \\
        -0.04(0.05)-0.09(0.2)-0.06(0.15) \\
        -0.09(0.05)+0(0.2)-0.08(0.15)
    \end{bmatrix} \\
    &=
    \begin{bmatrix}
        -0.0235 \\
        -0.0290 \\
        -0.0165
    \end{bmatrix}.
\end{aligned}
\]

\subsubsection*{3. Add the bias and apply ReLU}

\[
\begin{aligned}
    z_h
    &=
    \begin{bmatrix}
        -0.045 \\
        0.010 \\
        -0.075
    \end{bmatrix}
    +
    \begin{bmatrix}
        -0.0235 \\
        -0.0290 \\
        -0.0165
    \end{bmatrix}
    +
    \begin{bmatrix}
        0.1 \\
        0.1 \\
        0.1
    \end{bmatrix} \\
    &=
    \begin{bmatrix}
        0.0315 \\
        0.0810 \\
        0.0085
    \end{bmatrix}.
\end{aligned}
\]

Since all entries are positive,
\[
    h_t = \operatorname{ReLU}(z_h)
    =
    \begin{bmatrix}
        0.0315 \\
        0.0810 \\
        0.0085
    \end{bmatrix}.
\]

\subsubsection*{4. Compute the output layer}

The output pre-activation is
\[
    z_y = W_{hy}h_t + b_y.
\]

\[
\begin{aligned}
    z_y
    &=
    \begin{bmatrix}
        0.06 & -0.01 & -0.18 \\
        -0.06 & 0.13 & 0.14
    \end{bmatrix}
    \begin{bmatrix}
        0.0315 \\
        0.0810 \\
        0.0085
    \end{bmatrix}
    +
    \begin{bmatrix}
        0.1 \\
        0.1
    \end{bmatrix} \\
    &=
    \begin{bmatrix}
        0.09955 \\
        0.10983
    \end{bmatrix}.
\end{aligned}
\]

Again, both entries are positive, so ReLU leaves them unchanged:
\[
    y_t = \operatorname{ReLU}(z_y)
    =
    \begin{bmatrix}
        0.09955 \\
        0.10983
    \end{bmatrix}.
\]

\[
    \boxed{
    y_t=
    \begin{bmatrix}
        0.09955 \\
        0.10983
    \end{bmatrix}
    }
\]

\subsection*{Part (b): Compute the \(\delta\) values}

Let the target be
\[
    t=
    \begin{bmatrix}
        0 \\
        1
    \end{bmatrix}.
\]

Use the same convention as earlier assignments:
\[
    \delta_j = \frac{\partial E}{\partial z_j}.
\]

Since every hidden and output pre-activation is positive, all ReLU derivatives are
\[
    \operatorname{ReLU}'(z)=1.
\]

\subsubsection*{1. Output-layer deltas}

For each output neuron,
\[
    \delta = (y-t)\odot \operatorname{ReLU}'(z_y).
\]
Therefore,
\[
    \delta_y
    =
    \begin{bmatrix}
        0.09955-0 \\
        0.10983-1
    \end{bmatrix}
    =
    \begin{bmatrix}
        0.09955 \\
        -0.89017
    \end{bmatrix}.
\]

So the two output-neuron deltas are
\[
    \boxed{
    \delta_6=0.09955,\qquad
    \delta_7=-0.89017
    }.
\]

\subsubsection*{2. Hidden-layer deltas}

The hidden deltas at time \(t\) are
\[
    \delta_h
    =
    \left(W_{hy}^T\delta_y\right)\odot \operatorname{ReLU}'(z_h).
\]

First compute \(W_{hy}^T\delta_y\):
\[
\begin{aligned}
    W_{hy}^T\delta_y
    &=
    \begin{bmatrix}
        0.06 & -0.06 \\
        -0.01 & 0.13 \\
        -0.18 & 0.14
    \end{bmatrix}
    \begin{bmatrix}
        0.09955 \\
        -0.89017
    \end{bmatrix} \\
    &=
    \begin{bmatrix}
        0.0593832 \\
        -0.1167176 \\
        -0.1425428
    \end{bmatrix}.
\end{aligned}
\]

Because every hidden ReLU derivative is \(1\),
\[
    \delta_h
    =
    \begin{bmatrix}
        0.0593832 \\
        -0.1167176 \\
        -0.1425428
    \end{bmatrix}.
\]

Thus the hidden-neuron deltas are
\[
    \boxed{
    \delta_3=0.05938,\qquad
    \delta_4=-0.11672,\qquad
    \delta_5=-0.14254
    }.
\]

\section*{Final Answer}

\[
    \boxed{
    y_t=
    \begin{bmatrix}
        0.09955 \\
        0.10983
    \end{bmatrix}
    }
\]

\[
    \boxed{
    \delta_3=0.05938,\quad
    \delta_4=-0.11672,\quad
    \delta_5=-0.14254,\quad
    \delta_6=0.09955,\quad
    \delta_7=-0.89017
    }
\]

\section*{Exam Notes}

\begin{itemize}
    \item This is a \emph{single-time-step} recurrent calculation: \(h_{t-1}\) is
    given, so part (a) is just the hidden-state update followed by the output-layer
    computation.
    \item The hidden update has the standard RNN form
    \[
        h_t=\phi(W_{xh}x_t+W_{hh}h_{t-1}+b_h).
    \]
    \item Since every pre-activation is positive, every ReLU derivative is \(1\),
    which makes the \(\delta\) calculations especially simple.
    \item For this question, the hidden deltas are computed from the output layer at
    the same time step. No future-time contribution is needed because only one
    target/output step is being evaluated here.
\end{itemize}

\end{document}
